% !TeX program = lualatex
% Compile twice: lualatex 108.tex
% All references and prose are embedded; no BibTeX database or external inputs.
\documentclass[11pt,a4paper]{article}
\usepackage[margin=24mm,headheight=14pt]{geometry}
\usepackage{fontspec}
\setmainfont{Latin Modern Roman}
\setsansfont{Latin Modern Sans}
\setmonofont{Latin Modern Mono}
\usepackage{amsmath,amssymb,mathtools}
\usepackage{unicode-math}
\setmathfont{Latin Modern Math}
\usepackage{microtype}
\usepackage{enumitem,xurl,needspace}
\usepackage[dvipsnames]{xcolor}
\usepackage{hyperref}
\hypersetup{unicode=true,colorlinks=true,linkcolor=MidnightBlue,urlcolor=MidnightBlue,
 pdftitle={Research attempt: Problem 108},pdfauthor={Research notebook},bookmarksnumbered=true}
\usepackage{bookmark}
\usepackage{tocloft}
\setlength{\cftsecnumwidth}{3em}
\cftsetpnumwidth{2.5em}
\cftsetrmarg{4em}
\usepackage{titlesec}
\titleformat{\section}{\Large\bfseries\raggedright}{\thesection}{0.7em}{}
\usepackage{fancyhdr}
\pagestyle{fancy}\fancyhf{}
\fancyhead[L]{\small\sffamily Open Problems Math | Research notebook}
\fancyhead[R]{\small Problem 108 | 27 September 2026}
\fancyfoot[C]{\thepage}
\setlength{\parindent}{0pt}
\setlength{\parskip}{3pt plus 1pt}
\setlength{\emergencystretch}{3em}
\setcounter{tocdepth}{1}
\setcounter{secnumdepth}{1}
\setlist[itemize]{leftmargin=3em,itemsep=5pt,topsep=4pt,parsep=0pt}
\allowdisplaybreaks
\newcommand{\N}{\mathbb{N}}
\newcommand{\Z}{\mathbb{Z}}
\newcommand{\Q}{\mathbb{Q}}
\newcommand{\R}{\mathbb{R}}
\newcommand{\C}{\mathbb{C}}
\newcommand{\F}{\mathbb{F}}
\newcommand{\A}{\mathbb{A}}
\newcommand{\PP}{\mathbb P}
\newcommand{\E}{\mathbb{E}}
\newcommand{\eps}{\varepsilon}
\newcommand{\dd}{\,\mathrm d}
\newcommand{\sref}[2]{\hyperlink{source:#1:#2}{[S#2]}}
\newcommand{\eref}[2]{\hyperlink{extra:#1:#2}{[E#2]}}
% Turn the next switch off for a shorter reading copy. The quotations preserve
% every exactTarget field, including qualifications omitted from a short summary.
\newif\ifshowcatalogue
\showcataloguetrue
\newcommand{\cataloguescope}[1]{\ifshowcatalogue
 \paragraph{Exact scope in the supplied catalogue.}
 \begin{quote}\small #1\end{quote}\fi}
\usepackage{etoolbox}
\AtBeginEnvironment{itemize}{\small}
\numberwithin{equation}{section}
% Standalone export. Original section 108 preserved verbatim outside marked additions.
\begin{document}
\setcounter{section}{107}
% ================================================================
% problemId: problem.invariant-subspace-problem-for-hilbert-space-operators
% source releaseRank: 108; source releaseStatus: open
% exactTarget SHA-256: 032dd770aec1016bafc4bdebc7e1b550ce671688b112cf86ed76a571dbac92de
\clearpage
\section{\texorpdfstring{Invariant Subspace Problem for Hilbert Space Operators}{Invariant Subspace Problem for Hilbert Space Operators}}\label{108}
\subsection{Definitions and mathematical statement}
A complex Hilbert space is separable if it has a countable dense subset. A bounded operator $T:H\to H$ is linear and satisfies $\|Tx\|\le C\|x\|$. A closed linear subspace $M$ is invariant if $T(M)\subseteq M$. The assertion is
\[
 \forall T\in\mathcal B(H)\ \exists M\subset H:\qquad
 M\text{ closed},\quad\{0\}\ne M\ne H,\quad T(M)\subseteq M.
\]
The subspace need not be reducing: invariance under the adjoint $T^*$ is not required. The field is complex and $H$ is infinite-dimensional.
\par\smallskip\textit{Formulation and concept references:} \sref{108}{1}.
\cataloguescope{Does every bounded linear operator T on an infinite-dimensional separable complex Hilbert space H have a nontrivial closed invariant subspace, that is, a closed subspace M with \{0\} != M != H such that T(M) is a subset of M?}
\subsection{Short English statement}
Must every bounded operator on a separable complex Hilbert space preserve some proper nonzero closed subspace?
% BEGIN RESEARCH ADDITION 108
\subsection{Research attempt}

\paragraph{Current frontier.}
The 2025 survey in the supplied source distinguishes the unresolved Hilbert
space problem from negative results on other Banach spaces. It also records
the existence of almost-invariant half-spaces of defect at most one for
bounded operators on reflexive spaces: $TY\subseteq Y+F$ with
$\dim F\leq1$. The finite-dimensional error is not an invariant-subspace
theorem for $T$ itself. \sref{108}{1}

A 2026 preprint of Clemente--Gallardo-Guti\'errez proves, as its stated
Theorem 2.6, a hyperinvariant-subspace result for nonscalar \emph{hyponormal}
operators admitting a decomposition into a tridiagonal operator and a
trace-class operator. These additional hypotheses are retained here; neither
hyponormality nor such a decomposition is assumed for the catalogue's
arbitrary $T$. The recent preprint is not independently proof-audited in
this entry. \eref{108}{1}

\paragraph{Attack route.}
Every finite complex compression has invariant subspaces, but its invariant
vectors can escape to higher coordinates as the dimension grows. One route
would force convergence of the corresponding projections. A more economical
route is to preserve only two fixed vectors: one approximately inside the
subspaces and one approximately orthogonal to them. I prove that these two
anchors suffice, without needing convergence of the projections.

A second possible route starts with a defect-one almost-invariant subspace
and repairs it by an invariant graph. I also test that route; a concrete
nilpotent example shows why a repair over the \emph{same} initial half-space
cannot be automatic. Thus neither weak compactness nor rank-one defect is
silently treated as the missing theorem.

\paragraph{Attempt.}
\textbf{1. Finite sections and the exact selection problem.}
Choose an orthonormal basis $(e_j)_{j\geq1}$ of $H$ and let $P_n$ project onto
$H_n=\operatorname{span}\{e_1,\ldots,e_n\}$. Define
\[
 T_n=P_nTP_n,
\]
viewed as an operator on all of $H$, zero on $H_n^\perp$. Then
$\|T_n\|\leq\|T\|$ and $T_n\to T$ strongly. For each fixed $k\geq0$,
\begin{equation}\label{108:powers}
 T_n^k x\longrightarrow T^k x\qquad(x\in H).
\end{equation}
For $k\geq1$, use the telescoping identity
\[
 T_n^k-T^k=\sum_{j=0}^{k-1}T_n^{k-1-j}(T_n-T)T^j;
\]
each term converges strongly to zero and the outer powers are uniformly
bounded. No convergence uniform in $k$ is required or asserted.

Let $\mathcal I_n(T)$ be the set of all orthogonal projections $Q$ satisfying
\begin{equation}\label{108:finite-invariance}
 Q=P_nQ=QP_n,\qquad (P_n-Q)TQ=0.
\end{equation}
These are exactly the projections onto invariant subspaces of $T_n$ lying
in $H_n$. The set is nonempty, since it includes $0$ and $P_n$; Schur
triangularization also supplies nontrivial intermediate ranks when $n>1$.
As a subset of finite-dimensional matrices it is compact: the projection
and invariance equations are closed and $\|Q\|\leq1$.

For fixed unit vectors $a,b\in H$, set
\begin{equation}\label{108:objective}
 d_n(a,b)=\min_{Q\in\mathcal I_n(T)}
 \left(\|(I-Q)a\|^2+\|Qb\|^2\right).
\end{equation}
The minimum is attained by finite-dimensional compactness. This only proves
that each optimization problem is well posed; it does not say its minima
tend to zero.

\medskip
\textbf{2. The two-anchor lemma.}
Suppose there are fixed unit vectors $a,b$, an increasing sequence $n_j$, and
$Q_j\in\mathcal I_{n_j}(T)$ such that
\begin{equation}\label{108:anchors}
 \eta_j=\|(I-Q_j)a\|\longrightarrow0,
 \qquad \zeta_j=\|Q_jb\|\longrightarrow0.
\end{equation}
Then $T$ has a nonzero proper closed invariant subspace.

\emph{Proof.} Put $S_j=T_{n_j}$. Invariance gives
$S_j^kQ_ja\in\operatorname{ran}Q_j$ for every $k\geq0$. Therefore
\begin{align}
 |\langle T^ka,b\rangle|
 &\leq \|(T^k-S_j^k)a\|
       +|\langle S_j^k(I-Q_j)a,b\rangle|
       +|\langle S_j^kQ_ja,Q_jb\rangle|\nonumber\\
 &\leq\|(T^k-S_j^k)a\|+\|T\|^k(\eta_j+\zeta_j).
 \label{108:anchor-estimate}
\end{align}
At $k=0$ use $\|T\|^0=1$. For every fixed $k$, the right side tends to zero
by \eqref{108:powers} and \eqref{108:anchors}. Hence
$\langle T^ka,b\rangle=0$ for all $k\geq0$. The space
\begin{equation}\label{108:cyclic-space}
 M=\overline{\operatorname{span}}\{a,Ta,T^2a,\ldots\}
\end{equation}
contains the nonzero vector $a$ and is contained in $b^\perp$, so it is
proper. It is closed by definition. Boundedness of $T$ carries the algebraic
span into itself and permits passage to its norm closure, proving
$T(M)\subseteq M$. No condition involving $T^*(M)$ has been imposed.
$\square$

In particular, the following concrete selection statement would imply the
exact catalogue problem:
\begin{quote}
For every bounded operator $T$ and every chosen orthonormal basis, there
exist unit vectors $a\perp b$ such that
\begin{equation}\label{108:missing}
 \liminf_{n\to\infty}d_n(a,b)=0.
\end{equation}
\end{quote}
Choosing a subsequence of minimizers then supplies \eqref{108:anchors}.
Statement \eqref{108:missing} is \emph{not proved}. It is presented as a
sufficient finite-section selection principle, not as an established
equivalence or a theorem about all bases. The insistence that $a,b$ be the
same vectors at every stage is essential.

\medskip
\textbf{3. Comparison with a stronger compactness route.}
Suppose, more strongly, $Q_j\to Q$ strongly, with $Q$ a nonzero proper
orthogonal projection. From \eqref{108:finite-invariance},
\[
 (I-Q_j)TQ_j=(I-P_{n_j})TQ_j.
\]
For each fixed vector, the right side tends to zero: first replace $Q_jx$
by $Qx$, then use $P_{n_j}\to I$ strongly. Strong convergence and uniform
boundedness allow passage through the products on the left, giving
$(I-Q)TQ=0$. Thus $\operatorname{ran}Q$ is invariant.

This argument is valid, but its compactness assumption is absent in general.
The two-anchor lemma avoids having to compactify the entire sequence of
projections. It only asks for enough information to construct the cyclic
space \eqref{108:cyclic-space} and a single nonzero orthogonal vector.
The unproved selection principle remains a substantial requirement, rather
than a technical consequence of finite-dimensional triangularization.

\medskip
\textbf{4. Testing the defect-one repair idea.}
Let $Y$ be an almost-invariant half-space supplied by the theorem recorded
in \eref{108}{2}. Relative to $H=Y\oplus Y^\perp$, write
\[
 T=\begin{pmatrix}A&B\\ C&D\end{pmatrix},\qquad\operatorname{rank}C\leq1.
\]
A bounded graph $\{(y,Xy):y\in Y\}$ is invariant exactly when
\begin{equation}\label{108:riccati}
 C+DX=XA+XBX.
\end{equation}
This follows by applying $T$ to $(y,Xy)$ and equating its second component
with $X$ applied to its first. Hence a proposed repair by a bounded graph
has a precise operator Riccati equation to solve.

Rank one alone does not solve it. Take infinite-dimensional $Y,Z$,
$H=Y\oplus Z$, and
\[
 A=D=B=0,\qquad C:Y\to Z\text{ nonzero of rank one}.
\]
Then $T^2=0$ and $Y$ is almost invariant with defect one, but
\eqref{108:riccati} reduces to $C=0$, an impossibility. Thus there is no
invariant bounded graph over this particular $Y$. Nevertheless $Z$ itself
is an invariant subspace of $T$. This example refutes the universal
\emph{same-half-space graph-repair lemma}, not the invariant subspace
conjecture. Additional choices of half-space or genuinely different
constructions would be needed.

\paragraph{Stress test.}
\textbf{Escape at finite sections.}
Let $S e_j=e_{j+1}$ be the unilateral forward shift. Its compression to $H_n$
is a single nilpotent Jordan block. Its invariant subspace of dimension $m$
is
\[
 \operatorname{span}\{e_{n-m+1},\ldots,e_n\}.
\]
For fixed $m$ its projection tends strongly to zero. The same is true for
$m=\lfloor n/2\rfloor$: both its dimension and codimension tend to infinity,
but every fixed basis vector eventually lies outside it. A balance of
dimensions therefore gives no nonzero limiting subspace.

There is also a successful choice in this test case. Taking $m=n-1$ gives
$Q_n$ onto $\operatorname{span}\{e_2,\ldots,e_n\}$. The anchors $a=e_2$,
$b=e_1$ satisfy \eqref{108:anchors} exactly for $n\geq2$, and the resulting
space is $\overline{\operatorname{span}}\{e_2,e_3,\ldots\}$. Notice that
\[
 \|(I-Q_n)SQ_n\|=1
\]
for every $n\geq2$, because $e_n$ is sent to $e_{n+1}$. Requiring the
\emph{operator norm} of the invariance defect to vanish would exclude this
successful sequence. The pointwise limiting argument above needs no such
requirement.

\textbf{Weak limits of projections need not be projections.}
For $T=I$, let $R_n$ project onto the span of the $n$ orthonormal vectors
\[
 2^{-1/2}(e_i+e_{n+i}),\qquad1\leq i\leq n.
\]
Every $R_n$ is an exactly invariant projection lying in $H_{2n}$, but
$R_n\to\tfrac12 I$ in the weak operator topology. To see this, for each fixed
$i$ and all large $n$,
$R_ne_i=(e_i+e_{n+i})/2$; extend convergence of matrix entries by the uniform
norm bound and density. The limit is not idempotent. For any fixed unit
vector $u$,
\[
 \|R_nu\|^2\to\tfrac12,
 \qquad\|(I-R_n)u\|^2\to\tfrac12.
\]
So merely keeping vectors away from zero and from one in these norms does
not provide the two anchors.

More generally, if projections $Q_j$ converge weakly to a positive
contraction $A$, then
\begin{equation}\label{108:variance}
 \lim_j\|Q_ju-Au\|^2
 =\langle Au,u\rangle-\|Au\|^2
 =\langle(A-A^2)u,u\rangle.
\end{equation}
Expand the square, use $\|Q_ju\|^2=\langle Q_ju,u\rangle$, and pass to
weak limits in the cross term to prove the identity. The defect $A-A^2$
quantifies exactly why weak compactness fails to give the required strong
convergence. For $A=I/2$ it is $I/4$, not zero. Multiplication in the
invariance equation cannot simply be passed through a weak limit.

\textbf{A moving orthogonal vector is not an anchor.}
For every $n$ one can choose a unit vector $b_n$ perpendicular to
$a,Ta,\ldots,T^na$. This finite-dimensional observation is true even when
$a$ is cyclic for all of $H$. For example, with $T=S$, $a=e_1$, choose
$b_n=e_{n+2}$; then $b_n$ converges weakly to zero while
$\overline{\operatorname{span}}\{S^ke_1:k\geq0\}=H$.
Thus the selection statement must allow a suitable choice of $a$, and must
produce a \emph{fixed nonzero} $b$. It cannot be replaced by finitely many
orthogonality tests with an escaping witness.

The finite matrix tests in the companion script reproduce the Jordan-block
and projection identities. The infinite-dimensional conclusions are proved
above, not inferred from a numerical truncation. No compactness of the unit
sphere in norm, no spectral gap, and no diagonalizability of an arbitrary
operator is assumed.

\paragraph{Outcome.}
\textbf{Outcome: Conditional reduction.}

The proved two-anchor lemma reduces a positive solution to the explicit
selection principle \eqref{108:missing}; the latter remains unproved.
The argument includes a direct norm estimate and constructs the desired
closed invariant cyclic subspace once the anchors exist. It does not
supply those anchors for a general operator.

The stress tests identify three distinct failures: finite-section subspaces
can escape despite balanced dimensions; weak limits can lose idempotence;
and a rank-one almost-invariance defect need not admit a bounded graph
repair over the chosen half-space. These are counterexamples to intermediate
lemmas only. No bounded operator on the stipulated Hilbert space has been
shown to lack invariant subspaces, and no complete proof of the universal
positive assertion is claimed. No priority claim is made for the elementary
selection lemma or examples.

\par\needspace{20\baselineskip}
% END RESEARCH ADDITION 108
\subsection{Sources}
\begin{itemize}\raggedright
\item[\textnormal{[S1]}]\hypertarget{source:108:1}{} I. Chalendar and J. R. Partington, "Recent perspectives on the Invariant Subspace Problem," arXiv:2507.21834 (29 Jul 2025).
\url{https://arxiv.org/pdf/2507.21834}.
{\small\textit{Catalogue formulation source.}}
% BEGIN REFERENCE ADDITION 108
\item[\textnormal{[E1]}]\hypertarget{extra:108:1}{} N. Clemente and
E. A. Gallardo-Guti\'errez, \emph{Hyperinvariant subspaces of hyponormal
operators: A constructive decomposition approach}, arXiv:2606.00865v1
(2026), Theorem 2.6; see also Corollary 2.4.
\url{https://arxiv.org/html/2606.00865v1}.
{\small\textit{Consulted 27 September 2026; recent preprint, with the hyponormality and trace-class-perturbation hypotheses retained.}}
\item[\textnormal{[E2]}]\hypertarget{extra:108:2}{} I. Chalendar and
J. R. Partington, \emph{Recent perspectives on the Invariant Subspace
Problem}, arXiv:2507.21834v1 (29 July 2025), Section 2,
Definition 2.1 and Theorem 2.2, citing A. I. Popov and A. Tcaciuc,
\emph{Every operator has almost-invariant subspaces},
Journal of Functional Analysis 265 (2013), 257--265.
\url{https://arxiv.org/html/2507.21834v1}.
{\small\textit{Exact locator supplement to [S1], consulted 27 September 2026. The theorem is used only for almost invariance, not exact invariance.}}
% END REFERENCE ADDITION 108
\end{itemize}

\end{document}
